\section{Simplification and the LoD ladder}
\label{sec:lod}

\subsection{What simplification means here}

The octree's leaf lattice is the finest representation the shape pass ever
produces.  A coarser representation is obtained by removing all children of
a node and treating its eight interpolated corners as the node's SDF.  The
node then becomes a leaf at a coarser level and its dual vertex is solved
from the coarser eight samples.  The mesh is still a valid dual mesh; it
simply has fewer, larger cells.

Simplification is not a decimation of an existing triangle mesh.  It is a
decision made during the shape recursion, before any triangle exists.  This
is why the GPU can rebuild a chunk at any rung without a separate
simplification pass: the tree already stores the chosen level.

\subsection{The error criterion}

For a candidate parent cube and its eight child results, the simplifier
tests each surface child.  For every corner of the child cube it computes
the parent field's trilinear value and compares it with the child's stored
corner value:
\begin{equation}
  e = \max_{i\in\text{surface children}}\;
      \max_{j=0}^{7}
      \left| \operatorname{interp}(\mathbf{d}_{\text{parent}},
      C_{i,j}) - d_{i,j} \right|.
\end{equation}
The parent may collapse only when $e \le \tau L$, where $L$ is the parent
edge and $\tau$ a user tolerance.  A flat surface has small error and
collapses; a cliff, a thin wall, or a sharp CSG seam has large error and
stays subdivided.

\begin{lstlisting}[caption={The simplifier's core checks (abridged).},label={lst:simplify}]
SimplificationResult Simplifier::simplify(cube, sdf, children, chunkCube) {
    // 1. material consistency (texturing mode):
    //    all surface children must share one brush and be leaves
    // 2. chunk-border guard:
    //    never simplify a direct child of a chunk root
    // 3. re-interpolation error:
    //    for each surface child corner, |parent interp - child value|
    //    must be <= cube.getLengthX() * distance
    return result;
}
\end{lstlisting}

\subsection{The material and chain guards}

Two non-geometric guards are essential:

\begin{enumerate}
\item \textbf{Material boundary.}  If two surface children have different
  brushes, collapsing them would force one material onto the other.  The
  simplifier refuses.  This keeps a grass/rock boundary sharp even when the
  geometric error is small.
\item \textbf{Simplification chain.}  A child may be collapsed only if it is
  itself a leaf.  If a child still has children, collapsing the parent
  would skip a level of the ladder and break the hierarchy's monotonicity.
\end{enumerate}

\subsection{The chunk-border guard}

A chunk root's direct children have length $\Delta/2$.  They are the
boundary cells between two chunks.  If one chunk simplified such a cell and
its neighbour did not, the two dual meshes would choose different edge
vertices at the same boundary.  The stitching algorithm can handle
different cell sizes in the \emph{interior} of the adaptive walk, but the
chunk boundary is special: the two chunks are flattened independently and
never share a node table.  The guard therefore forbids simplifying the
direct children of a chunk root.  Coarser simplification is allowed only
deeper in the subtree, where both sides of a boundary see the same
neighbourhood and the break collector can stitch.

\subsection{Level propagation}

Each node's stored level is the minimum over its surface children plus one:
\begin{equation}
  lod(v) =
  \begin{cases}
    1, & v \text{ is a leaf at the frontier},\\
    0, & v \text{ is not a surface node},\\
    \min_{c\in\text{children}(v)} lod(c) + 1, & \text{otherwise}.
  \end{cases}
\end{equation}
The renderer uses the level as the rung index.  The chunk's own level is
propagated the same way but only along the chunk ladder:
\begin{equation}
  chunkLod(v) =
  \begin{cases}
    1, & v \text{ is a chunk root},\\
    \min_{c} chunkLod(c)+1, & \text{otherwise}.
  \end{cases}
\end{equation}
The published entry stores the 0-based \code{lod} (the shader/indirect gate
subtracts one implicitly).  A node with \code{chunkLod>0} and
\code{resultType==Surface} gets a version bump and is published.

\subsection{Why every rung is published}

It is tempting to publish only the finest rung and let the renderer collapse
it.  That fails for two reasons:

\begin{itemize}
\item \textbf{The tree is sparse.}  A coarse region may have no fine leaf at
  all; the only representation is the coarse node's eight samples.
\item \textbf{The dual mesh is not a simple triangle subset.}  A coarse cell
  produces a different dual vertex and a different ring than its eight
  fine children.  Decimating the fine mesh is not equivalent to solving the
  coarse dual mesh.
\end{itemize}

Publishing every rung gives the GPU gate a complete ladder.  The gate then
selects exactly one rung per region, so the extra rungs cost memory but not
draw calls.  The memory is bounded by the number of chunk-bearing nodes
(about one per $(\Delta/2,\Delta]$ cell) and is preallocated.

\subsection{Chunk versions and rebuilds}

Every published chunk carries a version.  The scene layer records the
version it has already emitted; a publish with the same version is
\code{UpToDate} and is skipped.  An edit bumps the version of the affected
chunk roots and their ancestors.  The renderer's pending queue then
contains a new entry for each affected rung.  The old span remains in the
pool until the new span is published, so no frame ever sees a hole.

\subsection{Cost model}

Let $C$ be the number of surface cells at the finest lattice, $N$ the
number of nodes in the tree, and $R$ the number of ladder rungs.  The
simplifier is linear in the number of surface children and corner tests:
\begin{equation}
  T_{\text{simplify}} = \Oh(N R \, 8 \cdot 8)
\end{equation}
in the worst case, but in practice most nodes fail the first material or
chain guard and do not reach the corner loop.  The dominant cost of a
shape apply is the field evaluation, not the simplifier.  The cache makes
repeated corner evaluations across neighbouring cubes cheap.
